% ============================================================================
% SEC01.TEX
% Section 1.1: Pengertian Game dan Elemen-elemen Game
% ============================================================================

\section{Pengertian Game dan Elemen-elemen Game}

\begin{learningoutcomes}
\item Mendefinisikan game secara formal dan akademis
\item Mengidentifikasi elemen-elemen pembentuk game (Mechanics, Dynamics, Aesthetics)
\item Menganalisis hubungan antara mechanics, dynamics, dan aesthetics
\item Menerapkan konsep MDA framework dalam analisis game
\end{learningoutcomes}

\subsection{Definisi Game}

\begin{definisi}[Game]
Game adalah aktivitas terstruktur yang melibatkan satu atau lebih pemain, dengan tujuan yang jelas, aturan yang mengatur interaksi, dan umpan balik yang memberikan informasi tentang kemajuan menuju tujuan \cite{Hunicke2004}.
\end{definisi}

Menurut \cite{Salen2004}, game memiliki karakteristik utama:
\begin{enumerate}
\item \textbf{Tujuan}: Setiap game memiliki tujuan yang harus dicapai pemain
\item \textbf{Aturan}: Aturan yang membatasi bagaimana pemain dapat mencapai tujuan
\item \textbf{Umpan Balik}: Informasi yang diberikan kepada pemain tentang status mereka
\item \textbf{Partisipasi Sukarela}: Pemain memilih untuk bermain dan menerima aturan
\end{enumerate}

\subsection{MDA Framework}

MDA (Mechanics, Dynamics, Aesthetics) merupakan framework analisis yang dikembangkan oleh \cite{Hunicke2004} untuk memahami game.

\subsubsection{Mechanics}

\begin{definisi}[Mechanics]
Mechanics adalah aturan dan sistem yang mengatur bagaimana game berfungsi. Mechanics mendefinisikan apa yang dapat dilakukan pemain dan bagaimana sistem game merespons aksi tersebut \cite{Hunicke2004}.
\end{definisi}

Contoh mechanics:
\begin{itemize}
\item Sistem pergerakan karakter
\item Mekanik pertempuran
\item Sistem inventori
\item Mekanik crafting
\item Sistem leveling dan progression
\end{itemize}

\subsubsection{Dynamics}

\begin{definisi}[Dynamics]
Dynamics adalah perilaku yang muncul dari interaksi mechanics saat game dimainkan. Dynamics adalah hasil dari mechanics yang bekerja bersama \cite{Hunicke2004}.
\end{definisi}

Contoh dynamics:
\begin{itemize}
\item Tension yang muncul dari time limit
\item Kompetisi antar pemain
\item Eksplorasi dan discovery
\item Risk vs reward decisions
\item Emergent gameplay
\end{itemize}

\subsubsection{Aesthetics}

\begin{definisi}[Aesthetics]
Aesthetics adalah pengalaman emosional yang dirasakan pemain saat bermain game. Aesthetics adalah respons emosional terhadap dynamics \cite{Hunicke2004}.
\end{definisi}

Jenis-jenis aesthetics menurut MDA framework:
\begin{enumerate}
\item \textbf{Sensation}: Game sebagai sense-pleasure
\item \textbf{Fantasy}: Game sebagai make-believe
\item \textbf{Narrative}: Game sebagai drama
\item \textbf{Challenge}: Game sebagai obstacle course
\item \textbf{Fellowship}: Game sebagai social framework
\item \textbf{Discovery}: Game sebagai uncharted territory
\item \textbf{Expression}: Game sebagai self-discovery
\item \textbf{Submission}: Game sebagai pastime
\end{enumerate}

\subsection{Hubungan Mechanics-Dynamics-Aesthetics}

\begin{figure}[h]
\centering
\begin{tikzpicture}[node distance=2cm]
    \node [block] (mechanics) {Mechanics\\Aturan Sistem};
    \node [block, below=of mechanics] (dynamics) {Dynamics\\Perilaku yang Muncul};
    \node [block, below=of dynamics] (aesthetics) {Aesthetics\\Pengalaman Emosional};
    
    \draw [line, ->] (mechanics) -- node[right] {Menghasilkan} (dynamics);
    \draw [line, ->] (dynamics) -- node[right] {Menciptakan} (aesthetics);
    \draw [line, ->, bend right=30] (aesthetics) to node[below] {Memengaruhi Desain} (mechanics);
\end{tikzpicture}
\caption{Hubungan Mechanics-Dynamics-Aesthetics}
\label{fig:mda_relationship}
\end{figure}

\subsection{Contoh Analisis Game Menggunakan MDA}

\begin{contoh}[Analisis Game Platformer]
\textbf{Game}: Super Mario Bros

\textbf{Mechanics}:
\begin{itemize}
\item Jump mechanics dengan gravity
\item Collision detection dengan platform
\item Power-up system (mushroom, flower)
\item Enemy AI behavior
\item Score system
\end{itemize}

\textbf{Dynamics}:
\begin{itemize}
\item Timing jump untuk menghindari musuh
\item Eksplorasi untuk menemukan power-up
\item Risk-reward dalam mengambil coin
\item Tension dari time limit
\end{itemize}

\textbf{Aesthetics}:
\begin{itemize}
\item Challenge: Mengatasi obstacle
\item Discovery: Menemukan secret area
\item Sensation: Visual dan audio feedback
\item Fantasy: Menjadi karakter hero
\end{itemize}
\end{contoh}

\subsection{Implementasi dalam Desain Game}

Saat merancang game, desainer harus mempertimbangkan:
\begin{enumerate}
\item \textbf{Dari Mechanics ke Dynamics}: Bagaimana mechanics akan menciptakan dynamics yang diinginkan?
\item \textbf{Dari Dynamics ke Aesthetics}: Bagaimana dynamics akan menghasilkan pengalaman emosional yang diinginkan?
\item \textbf{Iterasi}: Uji dan sesuaikan mechanics untuk mencapai aesthetics yang diinginkan
\end{enumerate}

\begin{catatan}
MDA framework membantu desainer memahami bahwa perubahan pada mechanics akan mempengaruhi dynamics dan akhirnya aesthetics. Oleh karena itu, desain game adalah proses iteratif yang memerlukan testing dan refinement.
\end{catatan}

\subsection{Ringkasan Section}

\begin{itemize}
\item Game adalah aktivitas terstruktur dengan tujuan, aturan, dan umpan balik
\item MDA framework membagi game menjadi Mechanics, Dynamics, dan Aesthetics
\item Mechanics adalah aturan sistem, Dynamics adalah perilaku yang muncul, Aesthetics adalah pengalaman emosional
\item Hubungan ketiganya adalah iteratif: Mechanics menghasilkan Dynamics, Dynamics menciptakan Aesthetics, dan Aesthetics mempengaruhi desain Mechanics
\item MDA framework berguna untuk analisis dan desain game
\end{itemize}
